\documentclass[10pt,a4paper]{amsart}
\usepackage[margin=19mm]{geometry}
\usepackage{amsmath,amssymb,mathtools}
\usepackage[scr=rsfs,cal=euler]{mathalfa}
\usepackage{xfrac}
\usepackage[unicode,hidelinks,bookmarks=false]{hyperref}
\newcommand{\ie}{i.e.\ }
\newcommand{\cf}{cf.\ }
\newcommand{\resp}{respectively\ }
\newcommand{\Subsection}{\subsection}
\providecommand{\nobreakdash}{\nobreak\hbox{-}}
\setlength{\emergencystretch}{2em}
\multlinegap0pt
\usepackage{amssymb}
\usepackage[shortlabels]{enumitem}
\usepackage{xcolor}
\usepackage{stackengine}
\newtheorem{theorem}{Theorem}
\usepackage{cleveref}
\crefname{theorem}{Theorem}{Theorems}
\newcommand{\M}{\mathcal{M}}
\newcommand{\N}{\mathcal{N}}
\newcommand{\Poin}{\mathrm{PI}  }
\title{Generalized Poincar\'e inequality for quantum Markov semigroups}
\author{Marius Junge, Jia Wang}
\date{Source: arXiv:2601.06005v2 (CC BY 4.0)}
\begin{document}
\maketitle

\noindent\emph{Source and licence.} Generalized Poincar\'e inequality for quantum Markov semigroups. Version arXiv:2601.06005v2 (CC BY 4.0), distributed by arXiv under the Creative Commons Attribution 4.0 International licence, http://creativecommons.org/licenses/by/4.0/, as recorded in the arXiv licence metadata for this exact version and shown on the abstract page https://arxiv.org/abs/2601.06005v2. Full licence notice: \texttt{licenses/arxiv-2601.06005-CC-BY-4.0.txt}.\par\smallskip

\noindent\textit{Editorial scope.} Counted unit: the article's theorem main:tracial (source0.tex lines 141--147). The article's complete original proof of it is delivered with it (source0.tex lines 1088--1131, one \texttt{proof} environment of 44 lines, which the article places away from the statement). No mathematics is written, repaired or re-derived: the statement and the proof are the article's own lines, in the article's own notation, and no retained line is substituted or re-flowed.  No further source-local context is needed: every cross-reference of every retained line resolves inside the two delivered spans.  Nothing was omitted from any retained span, so the package carries no omission record.  No retained line contains a citation, so the wrapper carries no bibliography and no bibliography entry.  No retained line contains a figure environment, an image-inclusion command or drawing code, so no image is delivered and the package declares no asset dependency.  Editorial formatting only, in addition: an AMS \texttt{amsart} preamble that restates the article's own declarations and macros used by the retained lines, namely the environments \texttt{theorem} on separate counters, together with the standard packages those lines need.  The retained environments print in this document as Theorem 1 in order of appearance, whereas the article prints the same items with the numbering of its own full document.  The complete native source of the article as distributed in the arXiv e-print for this exact version is bundled unchanged as \texttt{sources/192381/source0.tex}.
\begin{theorem}[$\Poin(p,p)$]\label{main:tracial}
Let $(\M,\tau)$ be a tracial von Neumann algebra and $(T_t)_{t\ge0}=e^{-tL}$ a $\tau$--symmetric quantum Markov semigroup with fixed-point algebra $\N$ and conditional expectation $E:\M\to\N$.
If $\{T_t\}_{t\ge0}$ has a spectral gap $\alpha>0$, then for every self-adjoint $x$ in the domain of $L$ and $p=2$ or $p\ge3$, it satisfies the Poincar\'e inequality $\Poin(p,p)$ with constant $\tfrac{p}{\sqrt{2\alpha}}$, i.e.,
\[
\|x-E(x)\|_{L^p(\M,\tau)}  \le  \frac{p}{\sqrt{2\alpha}} \|\Gamma(x,x)\|^{1/2}_{L^{p/2}(\M,\tau)}.
\]
\end{theorem}

\begin{proof}
Since $x=\Re x+\mathrm i\Im x$ with $\Re x,\Im x\in\M_{sa}$, the triangle inequality in $L_p$ gives
\[
\|x-E(x)\|_p
\le
\|\Re x-E(\Re x)\|_p
+
\|\Im x-E(\Im x)\|_p.
\]
We apply \Cref{main:tracial} to $\Re x$ and $\Im x$,
\[
\|\Re x-E(\Re x)\|_p
+
\|\Im x-E(\Im x)\|_p
\le
\frac p{\sqrt{2\alpha}}
\bigl(
\|\Gamma(\Re x,\Re x)\|_{p/2}^{1/2}
+
\|\Gamma(\Im x,\Im x)\|_{p/2}^{1/2}
\bigr).
\]
Since
\[
\|u\|_{\Gamma,p}:=\|\Gamma(u,u)\|_{p/2}^{1/2}
\]
is a seminorm and
\[
\Re x=\frac{x+x^*}{2},
\qquad
\Im x=\frac{x-x^*}{2\mathrm i},
\]
its triangle inequality gives
\[
\|\Gamma(\Re x,\Re x)\|_{p/2}^{1/2}
+
\|\Gamma(\Im x,\Im x)\|_{p/2}^{1/2}
\le
\|\Gamma(x,x)\|_{p/2}^{1/2}
+
\|\Gamma(x^*,x^*)\|_{p/2}^{1/2}.
\]
Combining the above estimates gives the stated corollary.
\end{proof}

\end{document}
